Benefiting from comprehensive training data and large-scale model parameters, large language models (LLMs) are increasingly employed to develop autonomous agents capable of simulating human behavior~\cite {qin2023toolllm,shinn2024reflexion,schick2024toolformer}.
These language agents have been explored as human proxies in various fields of behavioral research, such as psychological and sociological surveys~\cite{argyle2023out,chuang2024wisdom,xie2024largelanguagemodelagents}, and opinion dynamics modeling~\cite{mou-etal-2024-unveiling,liu2024skepticism}.
These social science studies often assume that LLMs exhibit \textbf{social intelligence}, the ability to navigate complex and multifaceted social goals during interactions with others~\cite{zhou2024sotopiainteractiveevaluationsocial}, given their ability to produce human-like responses and behaviors in certain scenarios.
Social intelligence is essential for fostering effective and harmonious interactions among agents and between agents and humans~\cite{xu2023exploring}.
However, the question remains: can LLMs truly navigate the intricacies of such interactions and exhibit social intelligence?
As illustrated in Figure~\ref{fig:intro}, LLM-driven agents still struggle with complex social situations. For example, when facing multiple goals in an office conversation, agent Emily fails to complete all her goals and cannot guess what Jordan has in mind.

\begin{figure}
    \setlength{\abovecaptionskip}{-0.2cm}
    \setlength{\belowcaptionskip}{-0.5cm}
    \centering
    \includegraphics[width=\linewidth]{latex/figs/intro.pdf}
    \caption{An illustration of challenging yet realistic social scenarios --- a family gathering and an office conversation, where the characters are driven by ChatGPT. While the dialogue could flow smoothly, Emily is unable to achieve her goals during the family gathering and fails to deduce Jordan's thoughts in an office setting.}
    \label{fig:intro}
\end{figure}

Early research on the social intelligence of LLMs relies on \textbf{static and non-interactive} tests that assess commonsense reasoning about social situations~\cite{sap2019socialiqacommonsensereasoningsocial,Zadeh_2019_CVPR,shapira2023clever,wilf2023social}, failing to capture the dynamic nature of social interactions. Recently, dynamic and goal-driven benchmarks~\cite{zhou2024sotopiainteractiveevaluationsocial, wang2024objectivelybenchmarkingsocialintelligence} have emerged to study social intelligence in interactive environments.
Although they have made some significant progress, they still exhibit three main limitations:
(1) \textbf{lack of scenario diversity}: existing studies~\cite{wang2024objectivelybenchmarkingsocialintelligence,sabour2024emobench} build social scenarios manually in a \textit{top-down} manner, resulting in a narrow set of common scenarios and goals, such as persuasion and collaboration~\cite{li2023camel}.
However, real-world interactions feature a broader spectrum of social goals and situational dynamics.
(2) \textbf{insufficient scenario complexity}: current work~\cite{xie2024largelanguagemodelagents,zhou2024sotopiainteractiveevaluationsocial} often limits interactions to two participants, each pursuing a single goal.
This oversimplifies real-life social interactions, where multiple actors engage simultaneously, each with multiple goals, as shown in Figure~\ref{fig:intro}.
As a result, it remains the performance of LLMs in group dynamics insufficiently studied. 
(3) \textbf{single-perspective observation}: while existing benchmarks primarily evaluate the extent to which agents achieve explicit social goals~\cite{li2023metaagents,wang2024objectivelybenchmarkingsocialintelligence}, they have neglected the concealment of private information and the inference of others' information, which are also important aspects in social interaction.

To this end, we introduce \textbf{AgentSense}: Benchmarking \underline{S}ocial Intellig\underline{en}ce of Language \underline{Agent}s through Interactive \underline{S}c\underline{e}narios.
We get inspiration from Dramaturgical Theory~\cite{goffman1959presentation}, which conceptualizes social interaction as a theatrical performance in which individuals assume specific roles within various settings.
We adopt a \textbf{bottom-up} approach, extracting scenarios from massive scripts to ensure that scenarios and social goals are diverse and grounded in real life.
In AgentSense, we construct 245 scenario templates from scripts, mitigating data leakage and expanding them into 1,225 scenarios, which cover existence goals to growth goals following the ERG theory~\cite{ALDERFER1969142}.

We then situate LLM-driven agents in the constructed scenarios, where each participant has social goals and may also safeguard some private information.
In an interactive environment, we simulate multi-turn interactions and observe agents' social intelligence from two aspects:
(1) goal completion: have the agents successfully achieved their social goals?
(2) implicit reasoning: can the agents accurately deduce others' private information?
These abilities are measured through interviews and multiple-choice questions, with our proposed PSI metric assessing profile sensitivity.
We find that LLMs struggle with complex social scenarios, particularly with high-level growth goals.

Our contributions are as follows:
\begin{itemize}
    \item We introduce AgentSense, a benchmark built on social scenarios derived from scripts using a bottom-up approach, distinguishing it from previous work.
    It encompasses diverse and challenging social scenarios, enabling a thorough evaluation of LLMs' social intelligence.
    \item We evaluate social intelligence from multiple aspects, considering both goal completion and information reasoning as well as profile sensitivity of social intelligence, through interviews with the agents and third-party judges.
    \item Our experiments reveal that LLMs struggle with complex scenarios and high-level goals, and their social intelligence is affected by profiles, interaction partners, and the balance between goals and privacy protection.
\end{itemize}